\section{Discussion}
\label{sec:discussion}

\subsection{Overview}

This section revisits the three research questions posed in
Section~\ref{sec:rq}, reflects on the challenges encountered during the
build, and discusses the limitations of the lab in the context of a
professional SOC environment. The intent is to be honest about what the
project does and does not represent, so its value as a portfolio artefact
rests on defensible claims rather than on overstated equivalence.

\subsection{RQ1: Can a Credible SOC Environment Be Built with Only Free and Open-Source Tools?}

The evidence from the build supports a qualified yes. Wazuh, Sysmon, and
\texttt{auditd}, deployed on VirtualBox with two commodity endpoint
operating systems, produced a working SIEM environment capable of ingesting
detection-grade telemetry, applying a mature rule set out of the box, and
raising meaningful alerts within seconds of simulated attacker activity.
Every tool used is deployed in production by real organisations, and the
concepts practised in the lab --- agent-based collection, log parsing,
rule-based correlation, alert triage --- transfer directly to commercial
platforms.

The qualification is that "credible" here means credible for the purposes of
individual learning and portfolio demonstration, not for the purposes of
running a live security programme for a real organisation. Multi-tenant
scalability, high availability, long-term compliance retention, integration
with commercial threat intelligence feeds, and 24-hour analyst rotation are
all absent. What the lab does credibly represent is the daily operational
core of a SOC analyst's work: interpreting alerts, correlating events across
sources, classifying activity, and producing structured triage
documentation.

\subsection{RQ2: What Additional Configuration is Required Beyond Defaults?}

This turned out to be the most instructive question of the three. A default
installation of Wazuh, a default installation of Windows, and a default
installation of Ubuntu, connected together, produced far less useful
telemetry than a naive reading of the vendor documentation would suggest.
Three specific configuration additions were required to reach
detection-grade visibility:

\begin{enumerate}[label=\textbf{\arabic*.},leftmargin=*]
    \item On Windows, Sysmon had to be installed and its dedicated event
          channel (\texttt{Microsoft-Windows-Sysmon/Operational}) had to be
          explicitly forwarded by the Wazuh agent through an added
          \texttt{<localfile>} block in \texttt{ossec.conf}. Without this,
          the agent received only standard Windows Security event log
          entries, which do not include the process, network, and
          file-creation records that modern behavioural detection depends
          on.
    \item On Linux, \texttt{auditd} was not present in the default Ubuntu
          Desktop install. It had to be installed, and audit rules for the
          \texttt{execve} system call had to be added and loaded before the
          Wazuh agent could see command-execution events.
    \item Simulation 4 identified a third gap: the default Wazuh syscheck
          File Integrity Monitoring configuration on Windows runs on a
          twelve-hour scan cycle rather than in real-time. A modification
          to a security-sensitive file such as the hosts file will
          eventually be detected, but not with the latency a live SOC would
          need. Closing this gap requires adding
          \texttt{realtime="yes"} to the relevant
          \texttt{<directories>} entries.
\end{enumerate}

The recurring theme is that vendor defaults optimise for broad coverage with
low overhead, not for security-critical corner cases. A deployment that
relies only on defaults will have predictable blind spots, and the work of
detection engineering is largely the work of identifying and deliberately
closing those blind spots. This directly motivates the planned Project~2 on
detection engineering with MITRE ATT\&CK.

\subsection{RQ3: To What Extent Does This Lab Represent a Professional SOC?}

Realistically, the lab represents perhaps twenty per cent of a professional
SOC's operational surface, and does so along one axis: the individual
analyst's interaction with a SIEM. The remaining eighty per cent, honestly
enumerated, includes:

\begin{itemize}[leftmargin=*]
    \item \textbf{Scale.} A production SIEM ingests events at rates several
          orders of magnitude beyond what two lab endpoints produce. Alert
          fatigue, triage prioritisation, and tuning at scale are among
          the hardest and most important skills in a real SOC, and the lab
          cannot exercise them.
    \item \textbf{Team workflow.} There is no Tier~2 or Tier~3 escalation,
          no shift handover, no on-call rotation, no incident coordination
          across multiple analysts. The lab is a single-analyst
          environment.
    \item \textbf{Network detection.} All monitoring in the lab is
          host-based. A production SOC also relies on network-tier
          detection through IDS/IPS platforms such as Suricata or Zeek and
          on full packet capture. This is out of scope here and identified
          as future work in Section~\ref{sec:conclusions}.
    \item \textbf{Cloud and SaaS telemetry.} Modern SOCs increasingly work
          with cloud audit logs, SaaS activity, and identity provider
          events. The lab is entirely on-premises virtualised
          infrastructure.
    \item \textbf{Threat intelligence enrichment.} Alerts in the lab are
          not enriched with external context such as IP reputation,
          malware family attribution, or campaign intelligence. In a
          production environment this enrichment is a major driver of
          triage decisions.
    \item \textbf{Response tooling.} There is no EDR quarantine, no SOAR
          playbook automation, and no case management ticketing. Response
          actions in the lab were performed by hand.
\end{itemize}

Within its scope, however, the lab does honestly reflect one specific
piece of SOC work: what an analyst does in the first sixty seconds after
an alert appears. That is a small share of a SOC's total surface but a
significant share of a Tier~1 analyst's daily work, and it is the piece
that is hardest to practise without a working SIEM to practise on.

\subsection{Challenges Encountered}

Three challenges shaped the build in ways worth documenting for future
readers.

\textbf{Resource constraint on the host machine.} The 16~GB of RAM on the
host was tight for running three virtual machines concurrently under any
sustained load. In practice the two endpoints were alternated for heavy
testing while the Wazuh manager remained continuously running. This is a
legitimate constraint for a home lab of this class, but it would prevent
any experiment that required both endpoints active simultaneously (for
example, lateral movement simulations between them).

\textbf{Vulnerability-detector disk consumption.} During extended idle
periods between build sessions, the Wazuh manager's vulnerability-detector
cache expanded until it consumed the entire 25~GB manager disk. The
symptom was manager service failure with a critical configuration error
that traced back to \texttt{touch: No space left on device}. The fix was
to clear the \texttt{/var/ossec/queue/vd\_updater/} and
\texttt{/var/ossec/queue/vd/} directories, restart the indexer and
manager, and (in the longer term) either grow the manager disk or disable
the vulnerability-detector module in \texttt{ossec.conf}. This incident is
worth recording because it is the kind of operational maintenance a lab
owner learns only by encountering, and because it illustrates how a SIEM
component can fail silently between active use.

\textbf{Configuration file structure across versions.} Some steps in
public Wazuh guides refer to configuration paths or agent commands that
have changed between major versions. Working from the vendor's own
documentation for the specific version installed
(\texttt{4.14.5}) rather than from third-party tutorials generally
avoided this issue, and this experience reinforces a lesson from the
methodology: recording exact software versions at build time is not
optional.

\subsection{Honest Reflection on Portfolio Value}

The value of this project as a portfolio artefact is not that it makes the
author qualified to operate a production SOC on its own; it does not.
The value is that it provides evidence of three specific things which are
frequently claimed by graduate applicants but rarely demonstrated:

\begin{itemize}[leftmargin=*]
    \item \textbf{Hands-on SIEM operation}, rather than only textbook
          familiarity, evidenced by the deployed environment and the
          documented triage of real alerts.
    \item \textbf{Awareness of the gap between defaults and
          detection-grade configuration}, evidenced by the deliberate
          endpoint telemetry work and by the negative-result FIM
          simulation.
    \item \textbf{Structured technical documentation and analyst-style
          reporting}, evidenced by the incident reports and by this
          report itself, both of which are held to a standard that is
          suitable for handover to a colleague.
\end{itemize}

These are the specific competencies most difficult for a graduate applicant
to demonstrate from coursework alone, and they are the specific
competencies this project was designed to develop.
